%------------------------------------------------------------------------------
% 4.1 المعمارية العامة للنظام
% المالك: AI (مهمة جانت ٢.٢.١ — التسليم 27/10/2026)
% الحجم المستهدف: 3-4 صفحات
% الغرض: الانتقال من التحليل (الفصل الثالث) إلى القرارات التصميمية الكلية:
%   لماذا هذه البنية تحديداً، وكيف تستجيب للمتطلبات غير الوظيفية.
%
% عناصر مطلوبة:
%   [ ] 4.1.1 وصف النمط المعماري: وكلاء لغويون منسقون برسم بياني حالة
%       (LangGraph) فوق بنية ثلاثية الطبقات (واجهة Next.js / خدمات FastAPI /
%       تخزين مزدوج PostgreSQL + Qdrant)، مع دور المنسق المركزي (Supervisor)
%   [ ] 4.1.2 شكل مخطط المكونات UML (fig:component-diagram) — TikZ جديد في
%       figures/tikz/component-diagram.tex: المكونات = واجهة المحادثة، خادم
%       API، الوكلاء الخمسة، طبقة العروض الذكية، قاعدة ERP (قراءة فقط)،
%       Qdrant، قاعدة بيانات النظام، وحدة التدقيق — مع واجهات موفرة/مطلوبة.
%       يمتاز عن الشكل المفاهيمي في الفصل الثاني بأنه تصميمي بمستوى
%       المكونات ومنافذها
%   [ ] 4.1.3 جدول القرارات المعمارية (tab:arch-decisions): القرار / البديل
%       المرفوض / المبرر / المتطلب المخدوم — 6-8 قرارات (لماذا LangGraph لا
%       CrewAI أو AutoGen؟ لماذا قراءة فقط من ERP؟ لماذا Qdrant لا pgvector؟
%       لماذا نشر Docker Compose؟ ...)
%
% تتبع: يستجيب للمتطلبات غير الوظيفية (الأداء، القابلية للتوسع، الخصوصية،
%   قابلية الصيانة) بأرقامها من جدول المتطلبات غير الوظيفية في الفصل الثالث
%------------------------------------------------------------------------------
\section{المعمارية العامة للنظام}
\label{sec:architecture}
~

\subsection{النمط المعماري}
\label{subsec:arch-pattern}
~

\subsection{مخطط المكونات}
\label{subsec:arch-components}
~

\subsection{القرارات المعمارية الرئيسة}
\label{subsec:arch-decisions}
~
